\section*{Chapter \ref{ch:b2:atmospheric-thermodynamics} --- Atmospheric Thermodynamics}

\begin{solution}{exo:b2:atmospheric-thermodynamics:1}
(a) $e = e_s(\qty{278.15}{K}) = \qty{8.73}{hPa}$ and $e_s(\qty{283.15}{K}) = \qty{12.32}{hPa}$.
(b) $8.73/12.32 = \qty{70.9}{\percent}$. (c) $r = 0.622 \times 8.73/(850 - 8.73) =
\qty{6.46}{g/kg}$ and $q = r/(1 + r) = \qty{6.42}{g/kg}$.
\end{solution}

\begin{solution}{exo:b2:atmospheric-thermodynamics:2}
(a) $\mathrm{d}\ln e_s/\mathrm{d}T = 5419/T^2$: \qty{7.26}{\percent} per kelvin at
\qty{273.15}{K}, \qty{6.31}{\percent} at \qty{293.15}{K}, \qty{5.90}{\percent} at
\qty{303.15}{K}. (b) $e_s(\qty{20}{\celsius}) = \qty{23.66}{hPa}$ and $e_s(\qty{25}{\celsius})
= \qty{32.26}{hPa}$; ratio 1.363, against $1.063^5 = 1.357$. The rate is a derivative of
$\ln e_s$, so over \qty{5}{K} the factor is $\exp(5 \times 0.0620) = 1.363$, with 0.0620 the
rate at the mid-point \qty{22.5}{\celsius}; compounding \qty{6.3}{\percent} five times is close
but not exact. (c) $23.66/23.39 = 1.011$: the formula is \qty{1.1}{\percent} too high, because
it keeps the latent heat of the triple point, while the latent heat decreases as the temperature
rises; the formula's slope is too steep, and its error grows with the distance from
\qty{0}{\celsius}.
\end{solution}

\begin{solution}{exo:b2:atmospheric-thermodynamics:3}
(a) $125 \times (30 - 14) = \qty{2000}{m}$. (b) $30 - 9.8 \times 2.0 = \qty{10.4}{\celsius}$.
(c) $14 - 1.8 \times 2.0 = \qty{10.4}{\celsius}$: temperature and dew point meet, as they must at
the level of saturation; the rule is this calculation. (d) $125 \times (33 - 14) = \qty{2375}{m}$:
the bases rise as the afternoon warms the ground air.
\end{solution}

\begin{solution}{exo:b2:atmospheric-thermodynamics:4}
(a) The dew-point curve lies on the temperature curve from the ground to about \qty{908}{hPa}:
the lowest half-kilometre is saturated, or nearly. (b) Between about 908 and \qty{889}{hPa}
the temperature curve turns to the right, to warmer values, as the pressure decreases: the
temperature rises with height. (c) LCL at about \qty{964}{hPa}, a few tens of metres up; LFC at
about \qty{730}{hPa}; EL at about \qty{235}{hPa}. (d) At the top of the inversion the dew point
is \qty{11}{K} below the temperature, and \qty{20}{K} below it at \qty{850}{hPa}: warm, very dry
air lies on top of cool, moist air. Its relative humidity is low, so it did not acquire its
warmth by condensation; air that sinks warms by compression at the dry adiabatic rate while
keeping its \omterm{def:b2:atmospheric-thermodynamics:moist-air}{mixing ratio}, which would give such a layer.
\end{solution}

\begin{solution}{exo:b2:atmospheric-thermodynamics:5}
(a) $T_v = 303.15 \times (1 + 0.020/0.622)/1.020 = \qty{306.76}{K}$; the approximation gives
$303.15 \times (1 + 0.61 \times 0.020) = \qty{306.85}{K}$. (b) $\rho = p/(R_dT_v) = 10^5/(287.0
\times 306.76) = \qty{1.1358}{kg/m^3}$, against $10^5/(287.0 \times 303.15) =
\qty{1.1494}{kg/m^3}$ for dry air: \qty{1.2}{\percent} lighter. (c) By $T_v - T =
\qty{3.6}{K}$. (d) A rising parcel needs only to be lighter than its surroundings: over warm seas
the moisture of the air adds a buoyancy equivalent to several kelvin, as much as a sizable
excess of temperature.
\end{solution}

\begin{solution}{exo:b2:atmospheric-thermodynamics:6}
(a) $e_s$ = 6.11, 23.66 and \qty{43.54}{hPa}; $r_s = 0.622\,e_s/(1000 - e_s)$ = 3.83, 15.07 and
\qty{28.31}{g/kg}; $\Gamma_s$ = 6.48, 4.20 and \qty{3.44}{K/km}. (b) The dry parcel cools by
\qty{9.8}{K}, the saturated one by \qty{4.2}{K}: the saturated parcel is \qty{5.6}{K} warmer.
(c) Higher up the air is colder, $r_s$ is smaller, little vapour is left to condense, and
$\Gamma_s$ returns towards $\Gamma_d$.
\end{solution}

\begin{solution}{exo:b2:atmospheric-thermodynamics:7}
(a) $e = e_s(\qty{291.65}{K}) = \qty{21.52}{hPa}$, $r = 0.622 \times 21.52/(971.1 - 21.52) =
\qty{14.1}{g/kg}$; $\theta = 292.15 \times (1000/971.1)^{0.2857} = \qty{294.6}{K}$; $\theta_e =
294.6 \exp[2.501 \times 10^6 \times 0.0141/(1004.5 \times 291.5)] = 294.6 \times 1.128 =
\qty{332.3}{K}$. (b) $e_s(\qty{264.25}{K}) = \qty{3.13}{hPa}$, $r_s = 0.622 \times 3.13/(500 -
3.13) = \qty{3.92}{g/kg}$; $\theta = 264.25 \times 2^{0.2857} = \qty{322.1}{K}$; the parcel is
saturated, so $T_{\text{LCL}} = T$ and $\theta_e = 322.1 \exp[2.501 \times 10^6 \times
0.00392/(1004.5 \times 264.25)] = \qty{334.3}{K}$. (c) $\theta$ rose by \qty{27.5}{K}, warmed
by the latent heat of the \qty{10.2}{g/kg} of vapour that condensed; $\theta_e$ changed by
\qty{0.6}{\percent}, within the accuracy of its approximate formula: it is conserved.
\end{solution}

\begin{solution}{exo:b2:atmospheric-thermodynamics:8}
(a) $\text{CAPE} = 287.0 \times 4 \times \ln(700/300) = 287.0 \times 4 \times 0.847 =
\qty{973}{J/kg}$. (b) $\sqrt{2 \times 973} = \qty{44.1}{m/s}$. (c) Between 1 and
\qty{1500}{J/kg}: positive, but not large. (d) A band of constant horizontal width between two
curves; by \cref{prop:b2:atmospheric-thermodynamics:areas} its area is proportional to $\int
\Delta T\,\mathrm{d}\ln p = 4 \ln(700/300)$, whatever the slant of the isotherms.
\end{solution}

\begin{solution}{exo:b2:atmospheric-thermodynamics:9}
(a) $D_c = \sqrt{3 \times 0.61 \times 80^3/2.1}\,\unit{nm} = \qty{668}{nm} =
\qty{0.668}{\micro m}$; $s_c = 2A/(3D_c) = 2 \times 2.1/(3 \times 668) = 0.00210 =
\qty{0.210}{\percent}$. (b) Solving $s_c = \sqrt{4A^3/(27\kappa D_d^3)}$ for $D_d$: $D_d =
\bigl(4A^3/(27\kappa s_c^2)\bigr)^{1/3}$, \qty{63}{nm} at \qty{0.3}{\percent} and \qty{131}{nm} at
\qty{0.1}{\percent}. (c) Every particle between 63 and \qty{131}{nm} now activates as well: more
droplets, and, for the same liquid water, smaller ones.
\end{solution}

\begin{solution}{exo:b2:atmospheric-thermodynamics:10}
(a) $w = \sqrt{2 \times 130} = \qty{16.1}{m/s}$, impossible for air at rest near the ground.
(b) The potential temperature of the inversion top is $293.65 \times (1000/889.1)^{0.2857} =
\qty{303.7}{K}$; a surface parcel with this $\theta$ has $T = 303.7 \times (971.1/1000)^{0.2857}
= \qty{301.2}{K}$, \qty{28.0}{\celsius}. (c) $125 \times (28.0 - 18.5) = \qty{1190}{m}$ above the
ground, above the old inversion top. (d) A front, which lifts the air on its slope; the slow ascent
of the air ahead of a disturbance aloft, which cools the lid itself.
\end{solution}

\begin{solution}{exo:b2:atmospheric-thermodynamics:11}
(a) When $r_s \to 0$, the numerator tends to $g$ and the denominator to $c_p$. (b) Write
$L^2r_s/(R_vT^2) = c_p \times \bigl(Lr_s/(R_dT)\bigr) \times LR_d/(R_vTc_p)$, with
$LR_d/(R_vTc_p) = 5.1$ at \qty{303}{K}: when $Lr_s/(R_dT) \gg 1$, also $L^2r_s/(R_vT^2) \gg c_p$,
and $\Gamma_s \to g\,(Lr_s/(R_dT))/(L^2r_s/(R_vT^2)) = gR_vT/(LR_d) = 9.81 \times 461.5 \times
303/(2.501 \times 10^6 \times 287.0) = \qty{1.91}{K/km}$. (c) With $T_d = T$,
$\Gamma_{\text{dew}} = gR_vT/(LR_d)$, the same. In that limit the latent heat is so large that
the parcel can condense almost none of its vapour without warming far above its saturation
point: it keeps its \omterm{def:b2:atmospheric-thermodynamics:moist-air}{mixing ratio} and follows the line of constant \omterm{def:b2:atmospheric-thermodynamics:moist-air}{mixing ratio}, along which the
temperature falls at $\Gamma_{\text{dew}}$. (d) At \qty{30}{\celsius} and \qty{1000}{hPa},
$Lr_s/(R_dT) = 0.81$, not much greater than 1, and $L^2r_s/(R_vT^2) = 4180$, four times $c_p$:
$\Gamma_s = \qty{3.4}{K/km}$, still far from \qty{1.9}{K/km}. The limit is not reached.
\end{solution}

\begin{solution}{exo:b2:atmospheric-thermodynamics:12}
(a) $2.868/2.599 = 1.104$ and $1.257/1.032 = 1.218$: \qty{10.4}{\percent} and
\qty{21.8}{\percent}. (b) Droplets activate at a few tenths of a per cent of \omterm{def:b2:atmospheric-thermodynamics:microphysics}{supersaturation},
tens of times less. (c) The crystals grow fast by deposition; they take vapour from the air,
which falls below saturation over the droplets; the droplets evaporate and their vapour goes to
the ice (the Bergeron process), until the cloud is glaciated. (d) At \qty{0}{\celsius} the two
saturation \omterm{def:b2:atmospheric-thermodynamics:moist-air}{vapour pressures} are equal: they meet at the triple point, and the ratio tends to one
near it.
\end{solution}

\begin{solution}{pb:b2:atmospheric-thermodynamics:1}
\textbf{1.} $e = e_s(\qty{291.65}{K}) = \qty{21.5}{hPa}$.
\textbf{2.} $r = 0.622 \times 21.5/(971.1 - 21.5) = \qty{14.1}{g/kg}$.
\textbf{3.} $q = 14.1/1.0141 = \qty{13.9}{g/kg}$.
\textbf{4.} $e_s(\qty{292.15}{K}) = \qty{22.2}{hPa}$; $21.5/22.2 = \qty{97}{\percent}$.
\textbf{5.} $0.61 \times 0.0141 \times 292.15 = \qty{2.5}{K}$.
\textbf{6.} $125 \times 0.5 = \qty{63}{m}$; $p = 971.1 \exp[-9.81 \times 62.5/(287.0 \times 294.7)]
= \qty{964}{hPa}$.
\textbf{7.} $\theta = 292.15 \times (1000/971.1)^{0.2857} = \qty{294.6}{K}$.
\textbf{8.} $16.9 - 17.3 = \qty{-0.4}{K}$: colder.
\textbf{9.} $(20.5 - 16.4)/(1102 - 924) = 4.1/178$: the temperature rises by \qty{23}{K/km}, a
lapse rate of \qty{-23}{K/km}. An inversion: absolutely stable, a lid.
\textbf{10.} $15.4 - 20.5 = \qty{-5.1}{K}$.
\textbf{11.} $e_s(\qty{-1.8}{\celsius}) = \qty{5.4}{hPa}$ and $e_s(\qty{18.6}{\celsius}) =
\qty{21.6}{hPa}$: \qty{25}{\percent}. Warm and dry, it can be air that sank and warmed by
compression, or air heated over dry ground upwind; it is not the moist air below.
\textbf{12.} $18.6 + 18.4 = \qty{37.0}{K}$ over $5725 - 1489 = \qty{4236}{m}$, \qty{8.73}{K/km}: between $\Gamma_s$ and
$\Gamma_d$, conditionally unstable.
\textbf{13.} $-18.4 - (-8.9) = \qty{-9.5}{K}$: extreme instability ($-8$ to $-10$).
\textbf{14.} From the ground up: 0, $-0.4$, $-0.2$, $-5.1$, $-4.9$, $+1.6$, $+9.5$, $+10.0$,
$+5.6$, $+1.5$, $-9.5$~\unit{K}.
\textbf{15.} Between 850 ($-4.9$) and \qty{700}{hPa} ($+1.6$), at the fraction $4.9/6.5 = 0.754$:
$\ln p = \ln 850 + 0.754 \ln(700/850)$, $p = \qty{734}{hPa}$; $z = 1489 + 0.754 \times 1621 =
\qty{2711}{m}$, \qty{2.37}{km} above the ground.
\textbf{16.} Between 250 ($+1.5$) and \qty{200}{hPa} ($-9.5$), at the fraction $1.5/11.0 =
0.136$: $p = \qty{242.5}{hPa}$ and $z = 10560 + 0.136 \times 1446 = \qty{10757}{m}$, about
\qty{200}{m} below the reported tropopause (\qty{10954}{m}).
\textbf{17.} The segments (mean $\Delta T$ times $\ln$ of the pressure ratio): $0.80 \times 0.0474
= 0.038$; $5.55 \times 0.3365 = 1.867$; $9.75 \times 0.2231 = 2.176$; $7.80 \times 0.2877 =
2.244$; $3.55 \times 0.1823 = 0.647$; $0.75 \times 0.0305 = 0.023$; sum \qty{6.995}{K}.
$\text{CAPE} = 287.0 \times 6.995 = \qty{2010}{J/kg}$, to three figures.
\textbf{18.} Between 1500 and \qty{2500}{J/kg}: large.
\textbf{19.} $-0.2 \times 0.0486 - 0.3 \times 0.0189 - 2.65 \times 0.0207 - 5.0 \times 0.0450 -
2.45 \times 0.1467 = -0.655$~\unit{K}; $\text{CIN} = 287.0 \times (-0.655) = \qty{-188}{J/kg}$.
\textbf{20.} The table keeps only eleven levels: between them the trapezium rule replaces the
curves by straight segments in $\ln p$ and misses their bends. The \qty{1}{hPa} grid follows the curves. The \omterm{def:b2:atmospheric-thermodynamics:virtual}{virtual temperature}
adds the buoyancy of the parcel's water vapour, more abundant in the saturated parcel than in
the dry air around it: a larger CAPE, a smaller CIN.
\textbf{21.} $\sqrt{2 \times 2010} = \qty{63}{m/s}$.
\textbf{22.} $10757 - 2711 = \qty{8046}{m}$ at $63.4/2 = \qty{31.7}{m/s}$: \qty{254}{s}, about four
minutes.
\textbf{23.} The weight of the condensed water the updraft carries; the mixing in of drier air
around it (entrainment), which cools it as droplets evaporate.
\textbf{24.} Parcel theory gives the updraft of a storm fed by this air a speed of about
\textbf{\qty[text-series-to-math=true]{63}{m.s^{-1}}}, an upper bound.
\end{solution}
